\section{Introduction}
\label{sec:intro}
Homomorphic encryption allows a service to evaluate a public computation without learning the encrypted inputs. Training is particularly demanding: the model changes at each step, gradients require transposed products, and a long optimization trajectory repeatedly exhausts the available multiplicative depth. A useful packing scheme must therefore account for the entire update, not just inference.

Mihara et al.~\cite{mihara2020} showed that diagonal-packed weights support transposition through rotations rather than ciphertext multiplications. Their interactive training design used a client to restore ciphertexts between updates. ReBoot~\cite{reboot2026} subsequently demonstrated non-interactive CKKS neural-network training, combining packed arithmetic, local learning signals, and bootstrapping. Non-interactivity itself is therefore not a novelty claim of this work. Our question is whether diagonal packing remains competitive when batch reductions, server-side layout maintenance, optimizer state, and bootstrapping are all charged.

We develop \sys{}, an MIT-licensed Rust implementation, and study three questions. \textbf{RQ1} asks how many key switches a complete training update requires compared with a ReBoot-style layout and the released ReBoot baseline. \textbf{RQ2} asks how bootstrap placement interacts with depth and compact state. \textbf{RQ3} asks what CUDA execution and its constituent optimizations contribute to measured training cost. A secondary continuity experiment revisits the original Iris setting; it is an interactive demonstration, not evidence of non-interactive training performance.

Our contributions are: (i) a rectangular, batched diagonal layout with a rotation-only transpose and explicit reduction costs; (ii) a compact-state non-interactive update whose operation counts can be checked against structural predictions; and (iii) an implementation and evaluation protocol that separates cryptographic assumptions, counter-derived comparisons, measured timing, and missing runs. Hoisting, baby-step giant-step (BSGS), sparse-secret encapsulation (SSE), and bootstrapping are established techniques; the contribution lies in their training-specific composition and accounting, not in inventing these primitives.

\paragraph{Status of the evidence.}
This is an interim full draft. The generated tables retain pending rows for the registered full-setting encrypted experiment and GPU training ablation. The completed encrypted subset runs use a different input dimension and cannot replace those rows. The released baseline and our preset also have different secret distributions and security estimates. We consequently do not claim that the registered equal-security, accuracy-matched research goal has been met. Section~\ref{sec:limits} identifies what would change that conclusion.
